% ==============================================================================
% CAPITOLO 10: INTERVALLI DI CONFIDENZA E INFERENZA INTERVALLARE
% ==============================================================================
\chapter{Intervalli di Confidenza e Inferenza Intervallare}
\label{cap:intervalli_di_confidenza}

\begin{abstract}
Il presente capitolo sviluppa la teoria della stima intervallare, lo strumento d'elezione per quantificare esplicitamente il grado di incertezza campionaria associato a una misura statistica. Si chiarisce l'autentico significato epistemologico frequentista del livello di confidenza $1-\alpha$, demistificando i comuni errori interpretativi. Viene formalizzato il metodo universale della quantità pivotale in 4 fasi e presentata la tassonomia esaustiva degli intervalli di confidenza: per la media di popolazioni gaussiane (con varianza nota e incognita), per proporzioni campionarie bernoulliane (con dimensionamento conservativo del campione), per il confronto tra due medie indipendenti (modelli a varianze note, omoschedastiche con pooled variance e modello di Welch) e per la varianza mediante pivot Chi-quadro.
\end{abstract}

% ==============================================================================
% SEZIONE 10.1: IL CONCETTO DI INTERVALLO DI CONFIDENZA
% ==============================================================================
\section{Il Concetto di Intervallo e l'Interpretazione Frequentista}
\label{sec:concetto_ic_frequentista}

Una stima puntuale come $\overline{x} = 10.2\text{ mm}$ fornisce un unico valore numerico, ma non comunica al decisore l'affidabilità o l'errore potenziale della misura: un conto è aver calcolato $\overline{x} = 10.2$ su $n=5$ pezzi ad alta variabilità, un altro è averlo ricavato da $n=10000$ provini di precisione. L'\textbf{inferenza intervallare} risolve questa lacuna associando alla stima una regione di valori plausibili accoppiata a un livello di affidabilità probabilistica prefissato $1-\alpha$ (tipicamente $95\%$ o $99\%$).

\begin{definizione}{Intervallo di Confidenza}{def_ic}
Sia $\mathbf{X} = (X_1, \dots, X_n)$ un campione casuale da una popolazione con parametro incognito $\theta$. Fissato un livello di significatività $\alpha \in (0, 1)$ (con coefficiente di confidenza $1-\alpha$), un \textbf{intervallo di confidenza al $(1-\alpha) \cdot 100\%$} per $\theta$ è una coppia di statistiche campionarie $L(\mathbf{X})$ e $U(\mathbf{X})$, con $L(\mathbf{X}) \le U(\mathbf{X})$, tali che:
\begin{equation}
    \mathbb{P}\left(L(\mathbf{X}) \le \theta \le U(\mathbf{X})\right) = 1 - \alpha
\end{equation}
\end{definizione}

\begin{figure}[htbp]
    \centering
    \begin{tikzpicture}[scale=0.8, >=Stealth]
    % Linea del vero parametro mu
    \draw[ultra thick, dashed, crimsonred] (0,-0.5) -- (0,10.5);
    \node[above, font=\bfseries\small, text=crimsonred] at (0,10.5) {Vero Parametro $\mu$};

    % Simulazione di 10 intervalli di confidenza estratti
    \foreach \y/\left/\right/\center/\captures in {
        10/-1.2/0.8/-0.2/1,
        9/-0.5/1.5/0.5/1,
        8/-1.8/0.2/-0.8/1,
        7/0.3/2.3/1.3/0,
        6/-1.0/1.0/0.0/1,
        5/-0.7/1.3/0.3/1,
        4/-1.5/0.5/-0.5/1,
        3/-0.2/1.8/0.8/1,
        2/-1.1/0.9/-0.1/1,
        1/-0.8/1.2/0.2/1
    } {
        \ifnum\captures=1
            \draw[thick, forestgreen] (\left, \y) -- (\right, \y);
            \draw[thick, forestgreen] (\left, \y-0.15) -- (\left, \y+0.15);
            \draw[thick, forestgreen] (\right, \y-0.15) -- (\right, \y+0.15);
            \fill[forestgreen] (\center, \y) circle (2pt);
        \else
            \draw[ultra thick, crimsonred] (\left, \y) -- (\right, \y);
            \draw[thick, crimsonred] (\left, \y-0.2) -- (\left, \y+0.2);
            \draw[thick, crimsonred] (\right, \y-0.2) -- (\right, \y+0.2);
            \fill[crimsonred] (\center, \y) circle (2.5pt);
            \node[right, font=\tiny\bfseries, text=crimsonred] at (\right, \y) {Non cattura $\mu$};
        \fi
    }

    \node[font=\footnotesize\itshape, text=navyblue, align=center] at (0,-1.2) {
        Nel $95\%$ dei campioni estratti, l'intervallo casuale $[\overline{X}_n - d, \overline{X}_n + d]$ conterrà il parametro fisso $\mu$.
    };
\end{tikzpicture}

    \caption{Significato frequentista di un intervallo di confidenza al $95\%$. La linea verticale rossa rappresenta il parametro vero $\mu$, fisso ed immutabile. Ciascun segmento orizzontale rappresenta l'intervallo casuale realizzato su un diverso campione estratto. Su un gran numero di repliche, esattamente il $95\%$ degli intervalli (in blu) cattura $\mu$, mentre il restante $5\%$ (in rosso) lo manca.}
    \label{fig:ic_ripetuti}
\end{figure}

\begin{esame}[Trappola Epistemologica: Cosa Significa Realmente "Confidenza al 95\%"?]
Una volta che il campione è stato estratto e l'intervallo numerico realizzato è stato calcolato (ad esempio, $[12.4, \; 15.8]$):
\begin{itemize}
    \item \textbf{ERRORE GRAVE:} Affermare che \emph{"la probabilità che $\theta$ sia compreso tra 12.4 e 15.8 è del 95\%"}. Nel paradigma frequentista, $\theta$ non è una variabile aleatoria, ma una costante fisica fissa. Di conseguenza, $\theta$ o appartiene a $[12.4, 15.8]$ (probabilità 1) o non vi appartiene (probabilità 0).
    \item \textbf{INTERPRETAZIONE CORRETTA:} Il valore $95\%$ misura l'\textbf{affidabilità a lungo termine della procedura statistica di stima}. Significa che se ripetessimo l'esperimento infinite volte estraendo nuovi campioni della medesima taglia, il $95\%$ degli intervalli casuali così costruiti conterrebbe il parametro vero $\theta$.
\end{itemize}
\end{esame}

% ==============================================================================
% SEZIONE 10.2: IL METODO DELLA QUANTITÀ PIVOTALE
% ==============================================================================
\section{Il Metodo Generale della Quantità Pivotale}
\label{sec:metodo_pivotale}

La tecnica universale per costruire intervalli di confidenza esatti e asintotici si basa sulla nozione di quantità pivotale.

\begin{definizione}{Quantità Pivotale}{def_pivot}
Una funzione $Q(\mathbf{X}, \theta)$ del campione e del parametro $\theta$ si dice \textbf{quantità pivotale} se:
\begin{enumerate}
    \item Dipende esplicitamente dai dati $\mathbf{X}$ e dal parametro target $\theta$;
    \item La sua legge di probabilità è \textbf{completamente nota} e non dipende da $\theta$ né da altri parametri ignoti di disturbo.
\end{enumerate}
\end{definizione}

\begin{figure}[htbp]
    \centering
    \begin{tikzpicture}[scale=0.85, >=Stealth]
    \draw[->, thick, navyblue] (-4.2,0) -- (4.2,0) node[right, font=\footnotesize] {$z$};
    \draw[->, thick, navyblue] (0,-0.2) -- (0,3.8) node[above, font=\footnotesize] {$\phi(z)$};

    % Area di confidenza 1 - alpha (centrale)
    \fill[coolblue!30, domain=-1.96:1.96, variable=\x]
        (-1.96,0) -- plot ({\x}, {3.5*exp(-\x*\x/2)}) -- (1.96,0) -- cycle;

    % Code di non confidenza alpha/2
    \fill[crimsonred!35, domain=-3.8:-1.96, variable=\x]
        (-3.8,0) -- plot ({\x}, {3.5*exp(-\x*\x/2)}) -- (-1.96,0) -- cycle;
    \fill[crimsonred!35, domain=1.96:3.8, variable=\x]
        (1.96,0) -- plot ({\x}, {3.5*exp(-\x*\x/2)}) -- (3.8,0) -- cycle;

    \draw[domain=-3.8:3.8, smooth, variable=\x, very thick, navyblue]
        plot ({\x}, {3.5*exp(-\x*\x/2)});

    \draw[thick, dashed, crimsonred] (-1.96,0) -- (-1.96, {3.5*exp(-1.96*1.96/2)});
    \draw[thick, dashed, crimsonred] (1.96,0) -- (1.96, {3.5*exp(-1.96*1.96/2)});

    \node[below, font=\small\bfseries, text=crimsonred] at (-1.96,0) {$-z_{\alpha/2}$};
    \node[below, font=\small\bfseries, text=crimsonred] at (1.96,0) {$+z_{\alpha/2}$};

    \node[font=\bfseries\small, text=navyblue] at (0,1.2) {Livello di Confidenza $1 - \alpha$ ($95\%$)};
    \node[font=\tiny\bfseries, text=crimsonred] at (-2.8,0.5) {$\alpha/2 = 2.5\%$};
    \node[font=\tiny\bfseries, text=crimsonred] at (2.8,0.5) {$\alpha/2 = 2.5\%$};
\end{tikzpicture}

    \caption{Costruzione dell'intervallo di confidenza sulla densità della quantità pivotale standardizzata $\mathcal{N}(0,1)$. L'area centrale racchiude la probabilità $1-\alpha$, delimitata dai quantili simmetrici $-z_{\alpha/2}$ e $+z_{\alpha/2}$, lasciando una massa $\alpha/2$ su ciascuna delle due code estreme.}
    \label{fig:ic_pivot}
\end{figure}

\begin{metodo}[Procedura Operativa in 4 Fasi per Costruire un Intervallo di Confidenza]
\begin{enumerate}
    \item \textbf{Scelta del Pivot:} Individuare la statistica $Q(\mathbf{X}, \theta)$ a distribuzione nota (es. Gaussiana $Z$, Student $T$, Chi-quadro $W$).
    \item \textbf{Definizione della Probabilità Centrale:} Dalla tavola della distribuzione, individuare i quantili critici $q_1, q_2$ tali che:
    \begin{equation}
        \mathbb{P}(q_1 \le Q(\mathbf{X}, \theta) \le q_2) = 1 - \alpha
    \end{equation}
    \item \textbf{Inversione Algebrica della Doppia Disuguaglianza:} Isolare il parametro $\theta$ al centro:
    \begin{equation}
        q_1 \le Q(\mathbf{X}, \theta) \le q_2 \iff L(\mathbf{X}) \le \theta \le U(\mathbf{X})
    \end{equation}
    \item \textbf{Valutazione Numerica:} Sostituire le realizzazioni campionarie $(\overline{x}, s, n)$ ottenendo $[L(\mathbf{x}), U(\mathbf{x})]$.
\end{enumerate}
\end{metodo}

% ==============================================================================
% SEZIONE 10.3: TASSONOMIA DEGLI INTERVALLI DI CONFIDENZA
% ==============================================================================
\section{Tassonomia Completa degli Intervalli di Confidenza}
\label{sec:tassonomia_ic}

\subsection{1. Media di Popolazione con Varianza $\sigma^2$ Nota (Intervallo Z)}
Se la popolazione è normale (o $n \ge 30$ per il CLT) con $\sigma$ noto, il pivot è $Z = \frac{\overline{X} - \mu}{\sigma / \sqrt{n}} \sim \mathcal{N}(0, 1)$.
\begin{equation}
    \operatorname{IC}_{1-\alpha}(\mu) = \left[ \overline{x} - z_{\alpha/2}\frac{\sigma}{\sqrt{n}}, \quad \overline{x} + z_{\alpha/2}\frac{\sigma}{\sqrt{n}} \right]
\end{equation}
\paragraph{Dimensionamento Campionario Ottimale:} Per garantire un margine di errore massimo $E = z_{\alpha/2}\frac{\sigma}{\sqrt{n}}$, la numerosità minima richiesta è:
\begin{equation}
    n \ge \left( \frac{z_{\alpha/2} \cdot \sigma}{E} \right)^2
\end{equation}

\subsection{2. Media di Popolazione con Varianza $\sigma^2$ Incognita (Intervallo $t$ di Student)}
Se la popolazione è normale con $\sigma$ incognito, sostituendo la deviazione standard campionaria $s$, il pivot è $T = \frac{\overline{X} - \mu}{S / \sqrt{n}} \sim t(n-1)$:
\begin{equation}
    \operatorname{IC}_{1-\alpha}(\mu) = \left[ \overline{x} - t_{\alpha/2, \, n-1}\frac{s}{\sqrt{n}}, \quad \overline{x} + t_{\alpha/2, \, n-1}\frac{s}{\sqrt{n}} \right]
\end{equation}

\subsection{3. Proporzione Bernoulliana (Intervallo Asintotico di Wald)}
Per una proporzione $p \in (0, 1)$ con stimatore $\widehat{p} = \frac{X}{n}$ su $n$ prove ($n\widehat{p} \ge 5, n(1-\widehat{p}) \ge 5$):
\begin{equation}
    \operatorname{IC}_{1-\alpha}(p) = \left[ \widehat{p} - z_{\alpha/2}\sqrt{\frac{\widehat{p}(1-\widehat{p})}{n}}, \quad \widehat{p} + z_{\alpha/2}\sqrt{\frac{\widehat{p}(1-\widehat{p})}{n}} \right]
\end{equation}
\paragraph{Dimensionamento Conservativo:} Poiché $\max_{p} p(1-p) = 0.25$ in $p=0.5$, il dimensionamento prudenziale che garantisce l'errore $E$ per qualsiasi $p$ è:
\begin{equation}
    n \ge \frac{z_{\alpha/2}^2}{4 E^2}
\end{equation}

\subsection{4. Differenza tra Due Medie Indipendenti $\mu_1 - \mu_2$}
\begin{itemize}
    \item \textbf{Varianze note $\sigma_1^2, \sigma_2^2$:}
    \begin{equation}
        \operatorname{IC}_{1-\alpha}(\mu_1 - \mu_2) = (\overline{x}_1 - \overline{x}_2) \pm z_{\alpha/2}\sqrt{\frac{\sigma_1^2}{n_1} + \frac{\sigma_2^2}{n_2}}
    \end{equation}
    \item \textbf{Varianze incognite ma uguali ($\sigma_1^2 = \sigma_2^2 = \sigma^2$ - Varianza Congiunta Pooled):}
    \begin{equation}
        s_p^2 = \frac{(n_1 - 1)s_1^2 + (n_2 - 1)s_2^2}{n_1 + n_2 - 2} \implies \operatorname{IC}_{1-\alpha}(\mu_1 - \mu_2) = (\overline{x}_1 - \overline{x}_2) \pm t_{\alpha/2, \, n_1+n_2-2} \, s_p \sqrt{\frac{1}{n_1} + \frac{1}{n_2}}
    \end{equation}
\end{itemize}

\subsection{5. Varianza di Popolazione Gaussiana $\sigma^2$}
Il pivot è $W = \frac{(n-1)S^2}{\sigma^2} \sim \chi^2(n-1)$:
\begin{equation}
    \operatorname{IC}_{1-\alpha}(\sigma^2) = \left[ \frac{(n-1)s^2}{\chi^2_{\alpha/2, \, n-1}}, \quad \frac{(n-1)s^2}{\chi^2_{1-\alpha/2, \, n-1}} \right]
\end{equation}
